\chapter{دوال \(L\) لديريشليه}

\section{نطاق الفصل وحالته}

\noindent\textbf{حالة الفصل:} \textenglish{DRAFT}. بدأ هذا الفصل بعد
تثبيت بوابة الأدلة في
\texttt{docs/CHAPTER\_07\_EVIDENCE\_2026-07-19.md}. يبني الفصل عائلة
دوال \(L\) التي تفصل الفئات المختزلة بترديد \(q\)، ويصل البنية الجبرية
للمحارف بالبنية التحليلية للسلاسل والمنتجات الأويلرية. البراهين القصيرة
الأساسية مكتوبة هنا، أما المعادلة الوظيفية وعدم الانعدام عند \(s=1\)
فمقتبسان صراحة من المراجع القياسية.

لا يغطي الفصل المناطق الكمية الخالية من الأصفار، أو أصفار زيغل، أو
الحدود الفعالة لأصغر أولي في متتالية حسابية، أو العزوم ودون التحدب، أو
فرضية ريمان المعممة. هذه موضوعات لاحقة، ولا يجوز استنتاج أي منها من
المادة التأسيسية الحالية.

\subsection{المتطلبات السابقة}

يفترض الفصل:

\begin{itemize}
  \item الجمع الجزئي لأبيل والتقارب المحلي المنتظم من الفصلين الثاني
  والثالث.
  \item الدوال الضربية والالتفاف الديريشلي من الفصل الرابع.
  \item المنتجات الأويلرية والمشتقة اللوغاريتمية من الفصل الخامس.
  \item الاستمرار التحليلي والمعادلة الوظيفية لزيتا من الفصل السادس.
  \item الحقائق الأساسية عن الزمر الأبيلية المنتهية ودالة غاما.
\end{itemize}

\subsection{نواتج التعلم}

بعد إتمام الفصل ينبغي أن يستطيع القارئ:

\begin{enumerate}
  \item تعريف محرف ديريشليه وتمييز الرئيسي والأولي والمستحث.
  \item استعمال علاقتي التعامد لبناء مؤشر فئة حسابية مختزلة.
  \item اشتقاق جداء أويلر لـ\(L(s,\chi)\) وعدم انعدامه عندما
  \(\Re(s)>1\).
  \item إثبات استمرار المحرف غير الرئيسي إلى \(\Re(s)>0\) بالجمع الجزئي.
  \item حساب الموصل والعوامل الأويلرية المحذوفة في حالة الاستحثاث.
  \item إثبات هوية مجموع غاوس وقيمته المطلقة للمحرف الأولي.
  \item قراءة تطبيع المعادلة الوظيفية وعامل الجذر من دون خلط المحرف
  بالمحرف المرافق.
  \item تحديد الأصفار البديهية للمحرف الأولي وشرح تناظرات الأصفار.
  \item بيان موضع عدم الانعدام عند الواحد في برهان مبرهنة ديريشليه.
  \item فصل النتائج النوعية في هذا الفصل عن المسائل الكمية المؤجلة.
\end{enumerate}

\section{لماذا نحتاج عائلة من الدوال؟}

يثبت المنتج الأويلري لزيتا أنها تجمع جميع الأوليات في دالة واحدة. لكن
السؤال عن الأوليات
\[
p\equiv a\pmod q,
\qquad (a,q)=1,
\]
يحتاج مرشحات تفصل فئة \(a\) من بقية الفئات المختزلة. قدم ديريشليه هذه
المرشحات في بحثه سنة 1837؛ وتتوفر ترجمة إنجليزية موثقة للبحث الأصلي في
\cite{Dirichlet1837Translation}.

الفكرة البنيوية هي تحليل الدالة المؤشرة للفئة \(a\) إلى ترددات منتهية.
هذه الترددات هي محارف زمرة الوحدات
\[
G_q=(\mathbb Z/q\mathbb Z)^\times.
\]
وبعد التحليل الجبري يتحول كل تردد إلى سلسلة ديريشليه ومنتج أويلري خاص
به. هكذا تصبح مسألة توزيعية على الأعداد الأولية مسألة عن قيم وأصفار
عائلة من الدوال العقدية.

\section{محارف ديريشليه}

\begin{definition}[محرف ديريشليه]
محرف ديريشليه بترديد، أو مقياس، \(q\geq1\) هو محرف زمري
\[
\chi:G_q\longrightarrow\mathbb C^\times
\]
يمدد إلى دالة \(\chi:\mathbb Z\to\mathbb C\) بوضع
\[
\chi(n)=0\quad\text{إذا }(n,q)>1,
\]
وبالاعتماد على فئة \(n\bmod q\) إذا كان \((n,q)=1\).
\end{definition}

ينتج من التعريف أن \(\chi\) دورية بدور \(q\)، وأنها ضربية تمامًا:
\[
\chi(mn)=\chi(m)\chi(n).
\]
ولأن \(G_q\) منتهية، فإن كل قيمة غير صفرية لـ\(\chi\) جذر وحدة، ومن ثم
\[
|\chi(n)|=
\begin{cases}
1,&(n,q)=1,\\
0,&(n,q)>1.
\end{cases}
\]

\begin{definition}[المحرف الرئيسي والمرافق والتكافؤ]
المحرف الرئيسي بترديد \(q\) هو
\[
\chi_0^{(q)}(n)=
\begin{cases}
1,&(n,q)=1,\\
0,&(n,q)>1.
\end{cases}
\]
والمحرف المرافق هو \(\overline\chi(n)=\overline{\chi(n)}\). على الوحدات
يكون \(\overline\chi=\chi^{-1}\). ونقول إن \(\chi\) زوجي إذا
\(\chi(-1)=1\)، وفردي إذا \(\chi(-1)=-1\). نرمز إلى التكافؤ بـ
\[
\mathfrak a(\chi)=
\begin{cases}
0,&\chi(-1)=1,\\
1,&\chi(-1)=-1.
\end{cases}
\]
\end{definition}

\begin{example}[المحرف الحقيقي بترديد \(4\)]
يوجد بترديد \(4\) محرفان، الرئيسي و
\[
\chi_4(n)=
\begin{cases}
0,&2\mid n,\\
1,&n\equiv1\pmod4,\\
-1,&n\equiv3\pmod4.
\end{cases}
\]
والمحرف \(\chi_4\) فردي لأن \(\chi_4(-1)=-1\). يبين الجدول قيمهما على
فئات البواقي:
\[
\begin{array}{c|rrrr}
n\bmod4&0&1&2&3\\ \hline
\chi_0^{(4)}(n)&0&1&0&1\\
\chi_4(n)&0&1&0&-1
\end{array}
\]
\end{example}

\subsection{حقيقة ازدواج الزمر المنتهية}

إذا كانت \(G\) زمرة أبيلية منتهية، فإن مجموعة محارفها
\(\widehat G\) زمرة عدد عناصرها \(|G|\)، وتفصل نقاط \(G\). وللتذكير
بسبب ذلك، تكتب مبرهنة بنية الزمر الأبيلية المنتهية
\[
G\simeq C_{m_1}\times\cdots\times C_{m_r}.
\]
للدورة \(C_m\) بالضبط \(m\) محارف، تحددها صورة مولدها بين جذور الوحدة
من الرتبة \(m\). وبأخذ الجداء نحصل على \(|G|\) محرفًا، وإذا كان عنصر
ما غير محايد فثمة عامل دوري يرى مركبته غير المحايدة. راجع
\cite{Apostol1976,MontgomeryVaughan2006}.

\begin{theorem}[علاقتا تعامد المحارف]
\label{thm:dirichlet-character-orthogonality}
\resultid{ANT-THM-07-01}
\provedhere
إذا كان \(\chi,\psi\) محرفين بترديد \(q\)، فإن
\[
\sum_{a\bmod q}\chi(a)\overline{\psi(a)}
=
\begin{cases}
\varphi(q),&\chi=\psi,\\
0,&\chi\ne\psi.
\end{cases}
\]
وكذلك، لأي عددين صحيحين \(a,b\)،
\[
\sum_{\chi\bmod q}\chi(a)\overline{\chi(b)}
=
\begin{cases}
\varphi(q),&a\equiv b\pmod q\ \text{و}\ (ab,q)=1,\\
0,&\text{خلاف ذلك}.
\end{cases}
\]
\end{theorem}

\begin{proof}
في العلاقة الأولى يكون \(\chi\overline\psi\) محرفًا على \(G_q\).
إذا كان رئيسيًا فكل حد على الوحدات يساوي \(1\)، فيكون المجموع
\(\varphi(q)\). وإذا لم يكن رئيسيًا نختار \(u\in G_q\) بحيث
\((\chi\overline\psi)(u)\ne1\). تغيير المتغير \(a\mapsto ua\) لا يغير
مجموعة الوحدات، لكنه يضرب المجموع في
\((\chi\overline\psi)(u)\)، فلا بد أن يكون المجموع صفرًا.

في العلاقة الثانية يكون المجموع صفرًا مباشرة إذا كان \(a\) أو \(b\)
غير وحدة. وإذا كانا وحدتين، نضع \(g=ab^{-1}\in G_q\). يصبح المجموع
\(\sum_{\chi\in\widehat G_q}\chi(g)\). إذا كان \(g=1\) فالقيمة
\(|\widehat G_q|=\varphi(q)\). وإذا كان \(g\ne1\)، يفصل ازدواج الزمرة
النقاط، فنختار \(\psi(g)\ne1\)، ويبين تبديل \(\chi\mapsto\chi\psi\)
أن المجموع \(S\) يحقق \(S=\psi(g)S\)، ومن ثم يساوي صفرًا.
\end{proof}

\begin{corollary}[مرشح الفئة الحسابية]
\label{cor:dirichlet-character-class-projector}
\resultid{ANT-COR-07-02}
\provedhere
إذا كان \((a,q)=1\)، فلكل عدد صحيح \(n\):
\[
\mathbf 1_{n\equiv a\ (q)}
=
\frac1{\varphi(q)}
\sum_{\chi\bmod q}\overline{\chi(a)}\chi(n).
\]
\end{corollary}

إذا تحقق \(n\equiv a\pmod q\)، فهو وحدة تلقائيًا؛ لذلك تغطي الصيغة
جميع \(n\) من دون شرط إضافي في الطرف الأيسر.

\section{سلسلة \(L\) والجداء الأويلري}

\begin{definition}
إذا كان \(\chi\) محرفًا بترديد \(q\)، نعرف في \(\Re(s)>1\)
\[
L(s,\chi)=\sum_{n=1}^{\infty}\frac{\chi(n)}{n^s}.
\]
\end{definition}

التقارب مطلق لأن \(|\chi(n)|\leq1\). والدالة \(\chi\) ضربية تمامًا،
فتنطبق عليها بنية الفصل الخامس.

\begin{theorem}[الجداء الأويلري وعدم الانعدام]
\label{thm:dirichlet-l-euler-product}
\resultid{ANT-THM-07-02}
\provedhere
إذا \(\Re(s)>1\)، فإن
\[
L(s,\chi)
=
\prod_p\left(1-\frac{\chi(p)}{p^s}\right)^{-1},
\]
والحاصل متقارب مطلقًا ومحليًا بانتظام، و\(L(s,\chi)\ne0\) في هذا النصف.
\end{theorem}

\begin{proof}
يعطي المنتج الأويلري العام من الفصل الخامس
\[
\sum_{n\geq1}\frac{\chi(n)}{n^s}
=
\prod_p\sum_{k\geq0}\frac{\chi(p)^k}{p^{ks}}.
\]
كل عامل محلي متسلسلة هندسية، فيعطي الصيغة المعلنة؛ وإذا كان \(p\mid q\)
فإن \(\chi(p)=0\) والعامل يساوي \(1\).

وعلى كل نصف مستوى \(\Re(s)\geq1+\delta\)، تتقارب
\(\sum_p|\chi(p)|p^{-\Re(s)}\) بانتظام. لذلك يتقارب مجموع اللوغاريتمات
\[
\sum_p-\log(1-\chi(p)p^{-s})
\]
محليًا بانتظام بعد اختيار الفرع الذي يساوي صفرًا عند اللانهاية. أسه
هذا المجموع هولومورفي ولا ينعدم، وهو المنتج الأويلري، فتنتج دعوى عدم
الانعدام.
\end{proof}

وبتفاضل اللوغاريتم في مجال التقارب المطلق نحصل على
\[
-\frac{L'}{L}(s,\chi)
=
\sum_{n=1}^{\infty}\frac{\Lambda(n)\chi(n)}{n^s},
\qquad \Re(s)>1.
\]
ومن مرشح الفئة، إذا \((a,q)=1\)،
\[
\sum_{\substack{n\geq1\\n\equiv a\ (q)}}
\frac{\Lambda(n)}{n^s}
=
-\frac1{\varphi(q)}
\sum_{\chi\bmod q}
\overline{\chi(a)}\frac{L'}{L}(s,\chi).
\]
هذه هي الصيغة التي تجعل أصفار دوال \(L\) مؤثرة في توزيع الأوليات بين
الفئات الحسابية.

\section{المحرف الرئيسي وقطب الواحد}

\begin{proposition}[تحليل دالة المحرف الرئيسي]
\label{prop:principal-character-l-function}
\resultid{ANT-PROP-07-01}
\provedhere
للمحرف الرئيسي بترديد \(q\):
\[
L(s,\chi_0^{(q)})
=
\zeta(s)\prod_{p\mid q}(1-p^{-s}).
\]
ومن ثم يمتد ميرومورفيًا إلى المستوى، وله قطب بسيط وحيد عند \(s=1\)،
وباقيه
\[
\operatorname*{Res}_{s=1}L(s,\chi_0^{(q)})
=
\prod_{p\mid q}\left(1-\frac1p\right)
=
\frac{\varphi(q)}q.
\]
\end{proposition}

\begin{proof}
في جداء أويلر تحذف العوامل الموافقة للأوليات القاسمة لـ\(q\):
\[
\prod_{p\nmid q}(1-p^{-s})^{-1}
=
\prod_p(1-p^{-s})^{-1}\prod_{p\mid q}(1-p^{-s}).
\]
هذا يثبت الهوية في \(\Re(s)>1\)، ثم يثبتها مبدأ الهوية حيث يمتد الطرفان.
وبما أن زيتا ذات قطب بسيط وباق \(1\) عند الواحد، والعوامل المنتهية
هولومورفية وغير منعدمة هناك، نحصل على الباقي المعلن.
\end{proof}

إذن المحرف الرئيسي يحتفظ بالقطب الذي يحمل الكثافة الكلية للأوليات، في
حين ينبغي أن تكون دوال المحارف غير الرئيسية منتظمة عند الواحد. هذا
الفارق هو قلب حجة ديريشليه.

\section{الاستمرار الأول للمحرف غير الرئيسي}

إذا كان \(\chi\ne\chi_0^{(q)}\)، فإن علاقة التعامد الأولى تعطي
\[
\sum_{a=1}^{q}\chi(a)=0.
\]
وبسبب الدورية، إذا كتبنا
\[
A_\chi(x)=\sum_{n\leq x}\chi(n),
\]
فإن الكتل الكاملة ذات الطول \(q\) تتلاشى، ومن ثم
\[
|A_\chi(x)|\leq q.
\]

\begin{theorem}[الاستمرار بواسطة الجمع الجزئي]
\label{thm:nonprincipal-l-half-plane-continuation}
\resultid{ANT-THM-07-03}
\provedhere
إذا كان \(\chi\) غير رئيسي، فإن
\[
L(s,\chi)
=
s\int_1^\infty A_\chi(x)x^{-s-1}\,dx
\]
في \(\Re(s)>1\)، ويعرف الطرف الأيمن استمرارًا هولومورفيًا وحيدًا إلى
نصف المستوى \(\Re(s)>0\).
\end{theorem}

\begin{proof}
تعطي صيغة أبيل، لأي \(N\geq1\)،
\[
\sum_{n\leq N}\frac{\chi(n)}{n^s}
=
A_\chi(N)N^{-s}
+s\int_1^N A_\chi(x)x^{-s-1}\,dx.
\]
إذا \(\Re(s)>0\)، فإن الحد الطرفي يؤول إلى الصفر لأن \(A_\chi\) محدودة،
ويتقارب التكامل مطلقًا. وعلى كل مجموعة مدمجة داخل هذا النصف يوجد
\(\delta>0\) بحيث يهيمن على الذيل ثابت مضروب في
\(\int_N^\infty x^{-1-\delta}\,dx\). لذلك التقارب محلي منتظم، والتكامل
هولومورفي. وهو يوافق السلسلة عندما \(\Re(s)>1\)، فيكون استمرارها الوحيد.
\end{proof}

\begin{corollary}[تمامية دالة المحرف غير الرئيسي]
\label{cor:nonprincipal-dirichlet-l-entire}
\resultid{ANT-COR-07-03}
\citedresult{\cite{Apostol1976}}
إذا كان \(\chi\) غير رئيسي، فإن \(L(s,\chi)\) تمتد إلى دالة تامة على
\(\mathbb C\).
\end{corollary}

لإثبات النتيجة هناك طريق بواسطة دالة هورفيتز. ففي
\(\Re(s)>1\):
\[
L(s,\chi)
=
q^{-s}\sum_{a=1}^{q}\chi(a)
\zeta\!\left(s,\frac aq\right).
\]
كل دالة هورفيتز في الطرف الأيمن تمتد ميرومورفيًا بقطب بسيط عند \(1\)
وباق \(1\)، ومجموع البواقي هو \(\sum_a\chi(a)=0\). ومن ثم يكون
\(L(s,\chi)\) تامًا للمحرف غير الرئيسي. هذه الاستمرارية القياسية لدالة
هورفيتز موثقة في \cite{Apostol1976}؛ وسيعطي قسم المعادلة الوظيفية طريقًا
آخر للمحارف الأولية ثم المستحثة.

\section{الموصل والمحارف الأولية}

\begin{definition}
نقول إن محرفًا \(\chi\) بترديد \(q\) مستحث من محرف \(\chi^*\) بترديد
\(f\mid q\) إذا
\[
\chi(n)=\chi^*(n)\chi_0^{(q)}(n)
\qquad(n\in\mathbb Z).
\]
والمحرف أولي إذا لم يستحث من مقياس أصغر. وأصغر مقياس \(f\) يستحث منه
\(\chi\) يسمى موصل \(\chi\).
\end{definition}

ينبغي عدم الخلط بين \emph{الرئيسي} و\emph{الأولي}: الأول يصف قيم
المحرف، والثاني يصف أصغر مقياس حقيقي له. فالمحرف الرئيسي بترديد \(q>1\)
مستحث من المحرف الوحيد بترديد \(1\)، ولذلك ليس أوليًا.

\begin{theorem}[وجود المحرف الأولي المولد ووحدانيته]
\label{thm:primitive-character-conductor}
\resultid{ANT-THM-07-04}
\provedhere
لكل محرف \(\chi\bmod q\) موصل وحيد \(f\mid q\)، ومحرف أولي وحيد
\(\chi^*\bmod f\) يستحث منه. وفي \(\Re(s)>1\):
\[
L(s,\chi)
=
L(s,\chi^*)
\prod_{\substack{p\mid q\\p\nmid f}}
\left(1-\frac{\chi^*(p)}{p^s}\right).
\]
\end{theorem}

\begin{proof}
اكتب \(q=\prod_p p^{e_p}\). تعطي مبرهنة البواقي الصينية
\[
G_q\simeq\prod_{p\mid q}G_{p^{e_p}},
\]
ومن ثم يتحلل \(\chi\) بوحيدانية إلى محارف محلية \(\chi_p\). لكل \(p\)
نختار أصغر \(f_p\in\{0,1,\ldots,e_p\}\) بحيث يكون \(\chi_p\) ثابتًا
على نواة الاختزال من الوحدات بترديد \(p^{e_p}\) إلى الوحدات بترديد
\(p^{f_p}\)، مع اعتبار \(f_p=0\) للمحرف المحلي الرئيسي. عندئذ يهبط
\(\chi_p\) إلى محرف على المقياس \(p^{f_p}\).

ضع \(f=\prod_p p^{f_p}\)، واجمع المحارف الهابطة بمبرهنة البواقي
الصينية. نحصل على \(\chi^*\bmod f\) يستحث منه \(\chi\). واختيار كل
\(f_p\) في الحد الأدنى يمنع هبوط \(\chi^*\) إلى قاسم أصغر، فيكون أوليًا.
وأي موصل آخر يعطي الأسس الدنيا المحلية نفسها، فتنتج الوحدانية.

يبقى بيان علاقة دوال \(L\). تتطابق العوامل الأويلرية لـ\(\chi\) و
\(\chi^*\) إلا عند أولي يقسم \(q\) ولا يقسم \(f\): عامل \(\chi\) هناك
يساوي \(1\)، في حين أن عامل \(L(s,\chi^*)\) هو
\((1-\chi^*(p)p^{-s})^{-1}\). ضرب العوامل المحذوفة يعطي الصيغة.
\end{proof}

\begin{example}[محرف مستحث بترديد \(6\)]
ليكن \(\chi^*\) المحرف غير الرئيسي بترديد \(3\)، حيث
\(\chi^*(1)=1\) و\(\chi^*(2)=-1\). المحرف المستحث بترديد \(6\) له
موصل \(3\)، وتكون
\[
L(s,\chi)=L(s,\chi^*)\bigl(1+2^{-s}\bigr),
\]
لأن \(\chi^*(2)=-1\). العامل المنتهي ليس تفصيلًا شكليًا: هو الفرق بين
المقياس المكتوب والموصل التحليلي الحقيقي.
\end{example}

\section{مجاميع غاوس}

نكتب
\[
\mathrm e(x)=e^{2\pi i x},
\qquad
\tau(n,\chi)=\sum_{a\bmod q}\chi(a)\mathrm e(an/q),
\]
ونضع \(\tau(\chi)=\tau(1,\chi)\).

\begin{theorem}[هوية مجموع غاوس]
\label{thm:primitive-gauss-sum}
\resultid{ANT-THM-07-05}
\provedhere
إذا كان \(\chi\) أوليًا بترديد \(q\)، فلكل عدد صحيح \(n\):
\[
\tau(n,\chi)=\overline{\chi(n)}\tau(\chi).
\]
وعلى الخصوص
\[
|\tau(\chi)|=\sqrt q,
\qquad
\tau(\chi)\tau(\overline\chi)=\chi(-1)q.
\]
\end{theorem}

\begin{proof}
إذا كان \((n,q)=1\)، فإن تغيير المتغير \(b=an\) يعطي
\[
\tau(n,\chi)
=
\sum_{b\bmod q}\chi(bn^{-1})\mathrm e(b/q)
=
\overline{\chi(n)}\tau(\chi).
\]
إذا كان \((n,q)>1\)، ضع \(d=(n,q)\). أولية \(\chi\) تعني وجود وحدة
\(u\equiv1\pmod{q/d}\) تحقق \(\chi(u)\ne1\)؛ وإلا كان \(\chi\)
ثابتًا على نواة الاختزال إلى مقياس أصغر. وبما أن
\(n(u-1)/q\in\mathbb Z\)، فإن تغيير \(a\mapsto au\) يعطي
\[
\tau(n,\chi)=\chi(u)\tau(n,\chi),
\]
ومن ثم \(\tau(n,\chi)=0\)، وهو الطرف الأيمن أيضًا لأن \(\chi(n)=0\).

نطبق الآن تعامد الإضافيات:
\[
\begin{aligned}
\sum_{n\bmod q}|\tau(n,\chi)|^2
&=
\sum_{a,b\bmod q}\chi(a)\overline{\chi(b)}
\sum_{n\bmod q}\mathrm e\!\left(\frac{n(a-b)}q\right)\\
&=q\sum_{a\bmod q}|\chi(a)|^2
=q\varphi(q).
\end{aligned}
\]
ومن الهوية الأولى يساوي الطرف الأيسر
\(\varphi(q)|\tau(\chi)|^2\)، فتنتج \(|\tau(\chi)|=\sqrt q\).
وأخيرًا
\[
\overline{\tau(\chi)}
=
\sum_{a\bmod q}\overline{\chi(a)}\mathrm e(-a/q)
=
\chi(-1)\tau(\overline\chi).
\]
نضرب في \(\tau(\chi)\) ونستعمل القيمة المطلقة لنحصل على الهوية الأخيرة.
\end{proof}

\section{الدالة المكتملة والمعادلة الوظيفية}

إذا كان \(\chi\) أوليًا بترديد \(q>1\)، ضع
\[
\mathfrak a=\mathfrak a(\chi)\in\{0,1\},
\]
وعرف
\[
\Lambda(s,\chi)
=
\left(\frac q\pi\right)^{(s+\mathfrak a)/2}
\Gamma\!\left(\frac{s+\mathfrak a}{2}\right)L(s,\chi),
\]
و
\[
\varepsilon(\chi)
=
\frac{\tau(\chi)}{i^{\mathfrak a}\sqrt q}.
\]
تعطي مبرهنة مجموع غاوس \(|\varepsilon(\chi)|=1\).

\begin{theorem}[المعادلة الوظيفية لدالة ديريشليه المكتملة]
\label{thm:primitive-dirichlet-l-functional-equation}
\resultid{ANT-THM-07-06}
\citedresult{\cite{Davenport2000,MontgomeryVaughan2006,IwaniecKowalski2004}}
إذا كان \(\chi\) أوليًا بترديد \(q>1\)، فإن \(\Lambda(s,\chi)\) دالة
تامة وتحقق
\[
\Lambda(s,\chi)
=
\varepsilon(\chi)\Lambda(1-s,\overline\chi).
\]
\end{theorem}

المسار القياسي للبرهان يبني دالة ثيتا الملتوية
\[
\theta_\chi(t)
=
\sum_{n\in\mathbb Z}n^{\mathfrak a}\chi(n)
e^{-\pi n^2t/q},
\]
ويطبق جمع بواسون. تظهر هوية مجموع غاوس عند تحويل فورييه، ثم يعطي تحويل
ميلين عامل غاما والقوة \((q/\pi)^{(s+\mathfrak a)/2}\). دقة الثابت
\(i^{\mathfrak a}\) وظهور \(\overline\chi\) سببان كافيان لعدم اختصار
هذا الاستشهاد إلى عبارة غير موثقة. تعرض المراجع المذكورة البرهان الكامل
بتطبيع مكافئ.

إذا كان \(\chi\) غير رئيسي وغير أولي، فإن مبرهنة الموصل تكتبه بدلالة
محرف أولي غير رئيسي وعوامل منتهية تامة. لذلك يمتد \(L(s,\chi)\) إلى
دالة تامة أيضًا، لكن معادلته الوظيفية النظيفة تخص \(\chi^*\)، لا المقياس
غير الأدنى \(q\).

\section{عدم الانعدام عند الواحد}

\begin{theorem}[مبرهنة ديريشليه لعدم الانعدام]
\label{thm:dirichlet-l-nonvanishing-at-one}
\resultid{ANT-THM-07-07}
\citedresult{\cite{Apostol1976,Davenport2000,MontgomeryVaughan2006}}
إذا كان \(\chi\) محرفًا غير رئيسي، فإن
\[
L(1,\chi)\ne0.
\]
\end{theorem}

هذه النتيجة أعمق من تقارب السلسلة عند الواحد. فالجمع الجزئي يثبت وجود
\(L(1,\chi)\)، لكنه لا يمنع أن تكون قيمتها صفرًا. وتختلف البراهين
الكلاسيكية بين المحارف غير الحقيقية والحقيقية؛ والحالة الحقيقية هي موضع
الدقة الرئيس. لذلك تعتمد الموسوعة هنا النتيجة من الكتب القياسية بدل
إخفاء هذه الفجوة داخل عبارة «واضح».

\section{الأصفار البديهية والتناظرات}

\begin{corollary}[الأصفار البديهية للمحرف الأولي]
\label{cor:primitive-dirichlet-l-trivial-zeros}
\resultid{ANT-COR-07-01}
\provedhere
ليكن \(\chi\) أوليًا بترديد \(q>1\)، ولتكن
\(\mathfrak a=\mathfrak a(\chi)\). عندئذ لـ\(L(s,\chi)\) أصفار بسيطة
عند
\[
s=-\mathfrak a-2m,
\qquad m=0,1,2,\ldots.
\]
وتسمى الأصفار البديهية. أما إذا كان \(\rho\) صفرًا غير بديهي، فإن
\[
L(\overline\rho,\overline\chi)=0,
\qquad
L(1-\rho,\overline\chi)=0,
\qquad
L(1-\overline\rho,\chi)=0,
\]
وإذا كان \(\chi\) حقيقيًا، تقع التناظرات الأربعة داخل الدالة نفسها.
\end{corollary}

\begin{proof}
لدالة غاما أقطاب بسيطة عند الأعداد الصحيحة غير الموجبة. لذلك يكون
\(\Gamma((s+\mathfrak a)/2)\) ذا قطب بسيط عند القيم المعلنة. أما
\(\Lambda(s,\chi)\) فتامة. وإذا كان \(s_0=-\mathfrak a-2m<0\)، فإن
المعادلة الوظيفية تربط \(\Lambda(s_0,\chi)\) بقيمة عند
\(1-s_0>1\)، حيث لا تنعدم \(L\) بفضل جداء أويلر؛ ولذلك تكون قيمة
\(\Lambda\) غير صفرية، ويلزم أن يلغي \(L\) قطب غاما بصفر بسيط.

تبقى الحالة \(s_0=0\)، ولا تقع إلا عندما \(\mathfrak a=0\). تربطها
المعادلة الوظيفية بـ\(s=1\)، وعدم الانعدام هناك هو المبرهنة السابقة،
فينتج الصفر البسيط أيضًا.

وأخيرًا تحقق السلسلة ثم الاستمرار التحليلي
\[
\overline{L(s,\chi)}=L(\overline s,\overline\chi).
\]
نجمع هذه الهوية مع المعادلة الوظيفية فنحصل على التناظر المعلن. وإذا
\(\chi=\overline\chi\)، تحفظ المرافقة الدالة نفسها، فتظهر المماثلة حول
المحور الحقيقي.
\end{proof}

لا يقال للمحرف غير الأولي إن جميع أصفاره البسيطة السابقة هي وحدها
«الأصفار البديهية»؛ فقد تضيف العوامل الأويلرية المنتهية المحذوفة أصفارًا
أخرى. كما لا يثبت هذا الفصل منطقة كمية خالية من الأصفار قرب
\(\Re(s)=1\).

\section{مثال كامل: دالة بيتا لديريشليه}

للمحرف \(\chi_4\) نحصل على
\[
L(s,\chi_4)
=
1-\frac1{3^s}+\frac1{5^s}-\frac1{7^s}+\cdots
=:\beta(s).
\]
المحرف أولي وفردي، و
\[
\tau(\chi_4)
=
\chi_4(1)i+\chi_4(3)(-i)=2i.
\]
إذن
\[
\varepsilon(\chi_4)=\frac{2i}{i\sqrt4}=1,
\]
والدالة المكتملة
\[
\Lambda(s,\chi_4)
=
\left(\frac4\pi\right)^{(s+1)/2}
\Gamma\!\left(\frac{s+1}{2}\right)\beta(s)
\]
تحقق
\[
\Lambda(s,\chi_4)=\Lambda(1-s,\chi_4).
\]
أصفار بيتا البديهية هي الأعداد الصحيحة السالبة الفردية
\(-1,-3,-5,\ldots\). يوضح المثال أن التكافؤ يحدد إزاحة عامل غاما، وأن
مجموع غاوس يحدد عامل الجذر.

\section{مبرهنة ديريشليه في المتتاليات الحسابية}

\begin{theorem}[ديريشليه]
\label{thm:dirichlet-primes-arithmetic-progressions}
\resultid{ANT-THM-07-08}
\provedhere
إذا كان \(q\geq1\) و\((a,q)=1\)، فهناك عدد لا نهائي من الأوليات
\(p\) التي تحقق
\[
p\equiv a\pmod q.
\]
\end{theorem}

\begin{proof}
لـ\(\sigma>1\)، يسمح التقارب المطلق بأخذ لوغاريتم المنتج الأويلري:
\[
\log L(\sigma,\chi)
=
\sum_p\frac{\chi(p)}{p^\sigma}+O(1)
\qquad(\sigma\to1^+).
\]
فالحدود ذات القوى \(p^{-m\sigma}\)، حيث \(m\geq2\)، مجموعها محدود
بانتظام قرب \(1\). نضرب في \(\overline{\chi(a)}\)، ثم نجمع على المحارف
ونستعمل التعامد، فنحصل على
\[
\sum_{\substack{p\equiv a\ (q)}}\frac1{p^\sigma}
=
\frac1{\varphi(q)}
\sum_{\chi\bmod q}
\overline{\chi(a)}\log L(\sigma,\chi)+O(1).
\]
الأوليات القاسمة لـ\(q\) لا تحقق الفئة المختزلة، ولا تؤثر حدود منتهية
على \(O(1)\).

للمحرف الرئيسي، تعطي قضية تحليله وقطب زيتا
\[
\log L(\sigma,\chi_0^{(q)})
=
\log\frac1{\sigma-1}+O(1).
\]
ولكل محرف غير رئيسي، تكون \(L(s,\chi)\) هولومورفية وغير منعدمة عند
\(1\) بالمبرهنة السابقة، فيوجد لوغاريتم محلي هولومورفي وتبقى
\(\log L(\sigma,\chi)=O(1)\). لذلك
\[
\sum_{\substack{p\equiv a\ (q)}}\frac1{p^\sigma}
=
\frac1{\varphi(q)}\log\frac1{\sigma-1}+O(1)
\longrightarrow\infty.
\]
لو كان عدد الأوليات في الفئة منتهيًا لبقي الطرف الأيسر محدودًا، وهذا
تناقض.
\end{proof}

البرهان المكتوب هنا نوعي ويعتمد في عقدته الوحيدة غير المثبتة داخليًا على
نتيجة \(L(1,\chi)\ne0\) المقتبسة. لا يعطي هذا البرهان كثافة تقاربية، ولا
حدًا لأصغر أولي، ولا انتظامًا في \(q\). هذه النتائج من نطاق فصل الأوليات
في المتتاليات الحسابية.

\section{خريطة الاعتماد المنطقي}

يمكن تلخيص السلسلة البرهانية هكذا:
\[
\begin{gathered}
\text{محارف زمرة الوحدات}
\Longrightarrow
\text{التعامد ومرشح الفئات},\\
\text{الضربية التامة}
\Longrightarrow
\text{الجداء الأويلري وعدم الانعدام في }\Re(s)>1,\\
\text{الدورية ومتوسط الصفر}
\Longrightarrow
\text{الاستمرار إلى }\Re(s)>0,\\
\text{الأولية ومجموع غاوس}
\Longrightarrow
\text{المعادلة الوظيفية},\\
L(1,\chi)\ne0
\Longrightarrow
\text{عدد لا نهائي من الأوليات في كل فئة مختزلة}.
\end{gathered}
\]
السهم الأخير لا يجوز عكسه بلا حجة إضافية، كما أن المعادلة الوظيفية وحدها
لا تثبت عدم الانعدام عند الواحد.

\section{أخطاء شائعة وحدود الاستدلال}

\begin{enumerate}
  \item \textbf{حصر قيم المحرف في \(\pm1\):} توجد محارف عقدية قيمها جذور
  وحدة غير حقيقية، كما يحدث بترديد \(5\).
  \item \textbf{الخلط بين الرئيسي والأولي:} المحرف الرئيسي بترديد كبير
  مستحث من المقياس \(1\)، فلا يكون أوليًا.
  \item \textbf{جعل مجال السلسلة مجال الدالة:} السلسلة تعرف مباشرة في
  \(\Re(s)>1\)، أما الدالة غير الرئيسية فتمتد إلى المستوى كله.
  \item \textbf{كتابة معادلة وظيفية نظيفة لمحرف مستحث:} التطبيع الطبيعي
  يخص المحرف الأولي ذي الموصل \(f\)، لا المقياس العارض \(q\).
  \item \textbf{نسيان المحرف المرافق:} المعادلة تربط \(\chi\) بـ
  \(\overline\chi\)، ولا تعود إلى الدالة نفسها إلا في حالات خاصة.
  \item \textbf{استنتاج عدم الانعدام عند الواحد من التقارب:} وجود قيمة
  منتهية لا يمنع أن تساوي صفرًا.
  \item \textbf{ترقية مبرهنة ديريشليه النوعية إلى تقدير كمي:} اللانهاية
  لا تعطي وحدها صيغة تقاربية أو حدًا فعالًا لأصغر أولي.
  \item \textbf{إخفاء أصفار العوامل المحذوفة:} المحارف غير الأولية قد
  تملك أصفارًا مصدرها الجداء المنتهي، لا عامل غاما.
\end{enumerate}

\section{أسئلة تفسيرية}

\begin{enumerate}
  \item لماذا يعطي تعامد المحارف مرشحًا للفئة المختزلة، لكنه لا يرشح
  فئة غير أولية مع \(q\) بالطريقة نفسها؟
  \item أين يستعمل التقارب المطلق في إثبات الجداء الأويلري وأخذ
  اللوغاريتم؟
  \item لماذا يكفي متوسط المحرف الصفري لجعل مجاميعه الجزئية محدودة؟
  \item ما الفرق التحليلي بين المقياس والموصل؟
  \item أين تدخل أولية المحرف تحديدًا في برهان هوية مجموع غاوس؟
  \item لماذا لا تكفي المعادلة الوظيفية وحدها لإثبات
  \(L(1,\chi)\ne0\)؟
  \item ما الحد الذي يتجه إلى اللانهاية في برهان ديريشليه، وما الذي يمنع
  هذا من إعطاء كثافة الأوليات مباشرة؟
\end{enumerate}

\section{تمارين متدرجة}

\begin{enumerate}
  \item اكتب محارف \(G_5\) الأربعة باختيار مولد \(2\)، وبين أن اثنين
  منها غير حقيقيين.
  \item استعمل علاقة التعامد الثانية لاشتقاق مرشح الفئة
  \(n\equiv3\pmod4\).
  \item للمحرف غير الرئيسي \(\chi_3\bmod3\)، احسب المجاميع
  \(\sum_{n\leq x}\chi_3(n)\) في دورة كاملة، واستنتج تقارب
  \(L(1,\chi_3)\).
  \item أثبت مباشرة من الجداء الأويلري صيغة
  \(-L'/L\) في \(\Re(s)>1\).
  \item احسب باقي \(L(s,\chi_0^{(12)})\) عند \(s=1\).
  \item تحقق من مثال الاستحثاث بترديد \(6\) بمقارنة العوامل الأويلرية
  عند \(p=2,3\) وعند بقية الأوليات.
  \item أثبت أن \(\tau(\chi_4)=2i\)، ثم تحقق من
  \(\tau(\chi_4)\tau(\overline{\chi_4})=-4\).
  \item بين أن المحرف الأولي الزوجي يحقق \(L(0,\chi)=0\)، وحدد النتيجة
  التي تحتاج إليها لإثبات بساطة هذا الصفر.
  \item أثبت أن الفرق بين
  \(\log L(\sigma,\chi)\) و\(\sum_p\chi(p)p^{-\sigma}\) محدود بانتظام
  عندما \(1<\sigma\leq2\).
  \item اشرح لماذا يثبت تباعد
  \(\sum_{p\equiv a(q)}p^{-\sigma}\) عند \(\sigma\to1^+\) وجود عدد لا
  نهائي من الأوليات، لكنه لا يثبت وحده
  \(\pi(x;q,a)\sim\operatorname{Li}(x)/\varphi(q)\).
\end{enumerate}

\subsection{حلول مختارة}

\paragraph{حل التمرين 3.}
قيم \(\chi_3\) في الدورة \(1,2,3\) هي \(1,-1,0\)، فمجموع الدورة صفر.
لذلك لا تتجاوز المجاميع الجزئية في القيمة المطلقة \(1\). ويعطي الجمع
الجزئي تقارب
\[
L(1,\chi_3)=\sum_{n\geq1}\frac{\chi_3(n)}n.
\]
هذا يثبت وجود القيمة، لا عدم انعدامها.

\paragraph{حل التمرين 5.}
من القضية الخاصة بالمحرف الرئيسي:
\[
\operatorname*{Res}_{s=1}L(s,\chi_0^{(12)})
=
\frac{\varphi(12)}{12}=\frac13.
\]

\paragraph{حل التمرين 9.}
من متسلسلة اللوغاريتم:
\[
\log L(\sigma,\chi)-\sum_p\frac{\chi(p)}{p^\sigma}
=
\sum_p\sum_{m\geq2}\frac{\chi(p)^m}{m p^{m\sigma}}.
\]
وقيمتها المطلقة لا تتجاوز
\[
\sum_p\sum_{m\geq2}\frac1{p^{m\sigma}}
\leq
\sum_p\frac{1}{p^2(1-p^{-1})},
\]
والمجموع الأخير متقارب، فيعطي حدًا مستقلًا عن
\(1<\sigma\leq2\).

\section{مسار القراءة والتوثيق}

للبناء التعليمي للمحارف ومجاميع غاوس ودوال \(L\)، انظر
\cite{Apostol1976}. وللحجة الكلاسيكية لمبرهنة ديريشليه وعدم الانعدام عند
الواحد، انظر \cite{Davenport2000}. وللموصل والمحارف الأولية والتطبيع
الدقيق للمعادلة الوظيفية، انظر
\cite{MontgomeryVaughan2006,IwaniecKowalski2004}.

سجل البحث السابق للتأليف محفوظ في وثيقة بوابة الأدلة، ويبين عبارات البحث
في zbMATH Open وSemantic Scholar وarXiv وOpenAlex، وما استعمل وما استبعد
ولماذا. ولا يدخل أي سجل حديث من تلك الفهارس في اعتماد المبرهنات
الكلاسيكية لهذا الفصل.

\section*{خلاصة الفصل}

حولت محارف ديريشليه مسألة الفئات الحسابية إلى تحليل فورييه منته على
\(G_q\). وحولت دوال \(L\) هذا التحليل إلى منتجات أويلر، ثم ربطت مجاميع
غاوس والعامل الأرخميدي قيم الدالة على جانبي الخط الحرج. بقيت العقدة
النوعية الحاسمة هي عدم الانعدام عند الواحد؛ وبها ينتج وجود عدد لا نهائي
من الأوليات في كل متتالية حسابية مختزلة. أما القياسات الكمية لتوزيع تلك
الأوليات فتبدأ في الفصل المخصص لها، لا في هذه المقدمة التأسيسية.
